% ===== latex-recomposed:document =====
% flm: style file flmrecomposed.cfeaebff.sty
\begin{flmUEpackedwrap}{{one}{2}{3}{4}{5}{6}{7}{8}{9}{ten}}the body\end{flmUEpackedwrap}
%%
% ===== latex-recomposed:style flmrecomposed.cfeaebff.sty =====
\NeedsTeXFormat{LaTeX2e}
\ProvidesPackage{flmrecomposed.cfeaebff}
% --- packages ---
% flm.latex.pkg.newunicodechar
\RequirePackage{newunicodechar}
% --- unicode characters written as UTF-8 ---
\newunicodechar{—}{\textemdash}
% --- commands and environments the library defines for the document ---
% flm.latex.def.ifblank
% \flmIfBlankTF{<content>}{<then>}{<else>}: the branch a user definition's own
% "is this blank?" test selects — <then> where <content> holds nothing at
% all, <else> otherwise. Not a customization point: change your own
% definition, not this. It is long, so that an argument holding a paragraph
% break is tested rather than stopping the compilation, and expandable, so
% that the test never stops the expansion of a definition that has to expand.
%
% This block uses expl3, LaTeX's programming layer, which
% is part of the format and needs no package; its names read only between
% \ExplSyntaxOn and \ExplSyntaxOff. \tl_if_empty:nTF is true exactly when its
% argument holds no token at all, which is the condition FLM's own blank test
% asks: a single space is content, and a comment is gone before the argument
% exists.
\ExplSyntaxOn
\providecommand\flmIfBlankTF[3]{\tl_if_empty:nTF {#1} {#2} {#3}}
\ExplSyntaxOff
% flm.latex.def.userargs
% \flmArg{<n>}{<arguments>}, \flmIfArgTF{<n>}{<arguments>}{<then>}{<else>} and
% \flmIfArgBlankTF{<n>,...}{<arguments>}{<then>}{<else>}: reading one argument of
% a user definition that takes more than nine of them. Such a definition is
% called with all its arguments in one group, one pair of braces around each
% ( \flmUDbig{{a}{b}{c}...} ), because a macro body has nine parameters; its
% body reads part <n> of that group through these. \flmArg gives the argument's
% content, \flmIfArgTF asks whether the argument holds \BooleanTrue, and
% \flmIfArgBlankTF asks whether the listed arguments hold nothing at all between
% them. Not a customization point: change your own definition, not this. Long,
% so that an argument holding a paragraph break is read rather than stopping the
% compilation, and expandable, so that reading an argument never stops the
% expansion of a definition that has to expand.
%
% \tl_item:nn is the LaTeX kernel's own: it picks part <n>, counting from one,
% and strips the braces around it. A test never reads a part directly --
% \IfBooleanTF{\flmArg{2}{#1}} stops with "Invalid argument" and
% \flmIfBlankTF{\flmArg{2}{#1}} answers the wrong branch without a word, since
% both look at the accessor's own tokens -- so the two tests expand the accessor
% first with \exp_args:Ne, which needs the \expanded primitive: every TeX engine
% has carried it since 2019, so it is older than the LaTeX this file asks for.
% What the part itself holds is left untouched: a space in an argument stays
% content, and a command in one is not expanded early.
%
% \flmIfArgBlankTF asks its question through \flmIfBlankTF, the block above,
% rather than spelling the test again: what "blank" means is written in one
% place regardless of a definition's size. A definition taking more than nine
% arguments therefore always brings that block into this file along with this
% one, whether or not it asks a blank test.
\ExplSyntaxOn
\providecommand\flmArg[2]{\tl_item:nn {#2} {#1}}
\providecommand\flmIfArgTF[4]{\exp_args:Ne \IfBooleanTF {\tl_item:nn {#2} {#1}} {#3} {#4}}
\providecommand\flmIfArgBlankTF[4]{\exp_args:Ne \flmIfBlankTF {\clist_map_tokens:nn {#1} {\tl_item:nn {#2}}} {#3} {#4}}
\ExplSyntaxOff
% --- user definitions ---
% flmUEpackedwrap (the document's packedwrap): defined unless the name is already taken, so a definition of it made before this file is loaded wins; \RenewDocumentEnvironment after loading this file changes it
\ProvideDocumentEnvironment{flmUEpackedwrap}{+m +b}{(note \flmArg{1}{#1}: \emph{#2} — end \flmArg{10}{#1})}{}
\endinput
